%% Response to Reviewers
%% Ion--Ice Reactions in Astrochemical Models
%% ============================================================
\documentclass[12pt]{article}

\usepackage[margin=1in]{geometry}
\usepackage{xcolor}
\usepackage{mdframed}
\usepackage{titlesec}
\usepackage{parskip}
\usepackage{hyperref}
\usepackage{enumitem}
\usepackage{booktabs}
\usepackage{amsmath}

%% ---- Colour scheme ----------------------------------------
\definecolor{reviewerblue}{RGB}{30, 90, 160}
\definecolor{commentbg}{RGB}{235, 242, 252}
\definecolor{responsebg}{RGB}{245, 250, 245}
\definecolor{commentframe}{RGB}{30, 90, 160}
\definecolor{responseframe}{RGB}{60, 140, 60}

%% ---- Framed environments ----------------------------------
\newmdenv[
  backgroundcolor=commentbg,
  linecolor=commentframe,
  linewidth=1.5pt,
  leftmargin=0pt,
  rightmargin=0pt,
  innerleftmargin=10pt,
  innerrightmargin=10pt,
  innertopmargin=8pt,
  innerbottommargin=8pt,
  skipabove=6pt,
  skipbelow=6pt,
]{reviewerbox}

\newmdenv[
  backgroundcolor=responsebg,
  linecolor=responseframe,
  linewidth=1.5pt,
  leftmargin=0pt,
  rightmargin=0pt,
  innerleftmargin=10pt,
  innerrightmargin=10pt,
  innertopmargin=8pt,
  innerbottommargin=8pt,
  skipabove=6pt,
  skipbelow=6pt,
]{responsebox}

%% ---- Section formatting -----------------------------------
\titleformat{\section}{\large\bfseries\color{reviewerblue}}{}{0em}{}[\titlerule]
\titleformat{\subsection}{\normalsize\bfseries}{}{0em}{}
\titleformat{\subsubsection}{\normalsize\itshape}{}{0em}{}

%% ---- Convenience commands ---------------------------------
\newcommand{\reviewlabel}{\textbf{\textcolor{commentframe}{Reviewer Comment:}}}
\newcommand{\responselabel}{\textbf{\textcolor{responseframe}{Authors' Response:}}}

%% ============================================================
\begin{document}

\begin{center}
  {\LARGE \textbf{Response to Reviewers}} \\[0.5em]
  {\large Ion--Ice Reactions in Astrochemical Models} \\[0.5em]
  {\normalsize \today}
\end{center}

\bigskip

We thank both reviewers for their careful reading of our manuscript and for
their constructive and detailed comments. We believe that addressing these
points has substantially improved the clarity and scientific rigour of the
work. Below we respond to each comment in turn. Reviewer comments are shown
on a \textcolor{commentframe}{blue} background; our responses are shown on a
\textcolor{responseframe}{green} background. Changes to the manuscript are
described in the responses and are highlighted in the revised manuscript.

\bigskip
\hrule
\bigskip

%% ============================================================
\section{Reviewer 1}
%% ============================================================

\textit{General comment:} The reviewer noted that the manuscript addresses a
scientifically relevant and timely topic, but requires clarification of the
theoretical framework, improved methodological transparency, and a more
explicit quantitative treatment of the modelling approach.

%% ------------------------------------------------------------
\subsection{Comment~1.1 -- Clarity of central objective}

\subsubsection{Reviewer Comment}
\begin{reviewerbox}
\reviewlabel{} The manuscript would benefit from a clearer statement of its
central objective in the introduction. It is not fully clear how this
modelling treatment of ion--ice chemistry differs from previous astrochemical
models that already include cosmic-ray-induced processes.
\end{reviewerbox}

\subsubsection{Authors' Response}
\begin{responsebox}
\responselabel{} 

% TODO: Insert response here.
\end{responsebox}

%% ------------------------------------------------------------
\subsection{Comment~1.2 -- Statement of novelty}

\subsubsection{Reviewer Comment}
\begin{reviewerbox}
\reviewlabel{} It should be clarified whether the novelty lies in updated
cross sections, revised rate parametrisations, explicit ion--solid reactions,
or a new coupling between laboratory data and kinetic simulations. A concise
paragraph explicitly stating the innovation of this study would significantly
strengthen the manuscript.
\end{reviewerbox}

\subsubsection{Authors' Response}
\begin{responsebox}
\responselabel{}

% TODO: Insert response here.
\end{responsebox}

%% ------------------------------------------------------------
\subsection{Comment~1.3 -- Energy deposition mechanisms}

\subsubsection{Reviewer Comment}
\begin{reviewerbox}
\reviewlabel{} The manuscript refers to ion--ice reactions but does not
clearly distinguish between the physical mechanisms underlying energy
deposition. For clarity and accuracy, please specify: (a) whether the model
considers electronic stopping ($S_e$) and nuclear stopping ($S_n$) separately;
(b) whether the main focus is on fast heavy ions (MeV--GeV cosmic rays) or
slower ions; (c) how the energy deposited per unit distance ($\mathrm{d}E/\mathrm{d}x$)
is handled; and (d) whether SRIM-based stopping powers were used or values
were taken from other sources.
\end{reviewerbox}

\subsubsection{Authors' Response}
\begin{responsebox}
\responselabel{}

% TODO: Insert response here.
\end{responsebox}

%% ------------------------------------------------------------
\subsection{Comment~1.4 -- Assumptions regarding homogeneous processing}

\subsubsection{Reviewer Comment}
\begin{reviewerbox}
\reviewlabel{} If the modelling approach assumes homogeneous chemical
processing throughout the ice mantle, this assumption must be explicitly
stated and justified. In reality, energy deposition is spatially localised
along ion tracks, and this can affect radical densities and reaction
probabilities. In addition, please clarify whether secondary electron cascades
are implicitly included, whether sputtering yields are considered, and whether
compaction or structural modifications of the ice are modelled.
\end{reviewerbox}

\subsubsection{Authors' Response}
\begin{responsebox}
\responselabel{}

% TODO: Insert response here.
\end{responsebox}

%% ------------------------------------------------------------
\subsection{Comment~1.5 -- Governing equations}

\subsubsection{Reviewer Comment}
\begin{reviewerbox}
\reviewlabel{} The manuscript would be strengthened by explicitly presenting
the governing equations used in the astrochemical simulations. If the model
alters the standard rate equation formalism, it is crucial to clearly
illustrate the following aspects, including explicit definitions and additional
ion-induced reaction rates. If ion-induced processes are introduced as
additional destruction or formation terms, please provide the exact rate
expression. All variables must be defined upon first use, and units must be
consistent throughout. If laboratory cross sections are converted into
astrophysical reaction rates, the scaling procedure and underlying assumptions
should be presented quantitatively.
\end{reviewerbox}

\subsubsection{Authors' Response}
\begin{responsebox}
\responselabel{}

% TODO: Insert response here.
\end{responsebox}

%% ------------------------------------------------------------
\subsection{Comment~1.6 -- Laboratory data to astrophysical timescales}

\subsubsection{Reviewer Comment}
\begin{reviewerbox}
\reviewlabel{} The manuscript needs a more transparent explanation of how
laboratory irradiation data are translated into astrophysical timescales.
Please clarify: what ion flux is assumed (particles~cm$^{-2}$~s$^{-1}$);
whether a monoenergetic approximation is used; how the cosmic-ray energy
spectrum is treated; and whether attenuation inside the cloud is considered.
\end{reviewerbox}

\subsubsection{Authors' Response}
\begin{responsebox}
\responselabel{}

% TODO: Insert response here.
\end{responsebox}

%% ------------------------------------------------------------
\subsection{Comment~1.7 -- Incorporation into the chemical network}

\subsubsection{Reviewer Comment}
\begin{reviewerbox}
\reviewlabel{} The manuscript should clarify how ion--ice reactions are
incorporated into the baseline chemical network. If laboratory cross sections
are converted into astrophysical reaction rates, the scaling procedure and
underlying assumptions should be presented quantitatively.
\end{reviewerbox}

\subsubsection{Authors' Response}
\begin{responsebox}
\responselabel{}

% TODO: Insert response here.
\end{responsebox}

%% ------------------------------------------------------------
\subsection{Comment~1.8 -- Competing processes}

\subsubsection{Reviewer Comment}
\begin{reviewerbox}
\reviewlabel{} How are competing processes (e.g., UV photolysis, thermal
diffusion, grain surface reactions) treated relative to ion processing? A
schematic representation of the modified network or a summary table of new
ion-induced reactions would be useful.
\end{reviewerbox}

\subsubsection{Authors' Response}
\begin{responsebox}
\responselabel{}

% TODO: Insert response here.
\end{responsebox}

%% ------------------------------------------------------------
\subsection{Comment~1.9 -- Quantitative discussion of results}

\subsubsection{Reviewer Comment}
\begin{reviewerbox}
\reviewlabel{} The results section would benefit from a more quantitative
discussion. Please provide: direct comparisons between models with and without
ion--ice chemistry; percentage changes in key molecular abundances; and
identification of the species most strongly affected by ion processing.
\end{reviewerbox}

\subsubsection{Authors' Response}
\begin{responsebox}
\responselabel{}

% TODO: Insert response here.
\end{responsebox}

%% ------------------------------------------------------------
\subsection{Comment~1.10 -- Minor editorial points}

\subsubsection{Reviewer Comment}
\begin{reviewerbox}
\reviewlabel{} Please verify consistent notation throughout; ensure cross
sections are reported with correct units (cm$^2$); please check dimensional
consistency in all rate expressions; and ensure all references cited in the
discussion are appropriate and up to date.
\end{reviewerbox}

\subsubsection{Authors' Response}
\begin{responsebox}
\responselabel{}

% TODO: Insert response here.
\end{responsebox}

\bigskip
\hrule
\bigskip

%% ============================================================
\section{Reviewer 2}
%% ============================================================

\textit{General comment:} The reviewer noted that this work appears to be one
of the first of its kind, opening the door to more elaborate studies, but
raised several specific concerns regarding clarity and completeness.

%% ------------------------------------------------------------
\subsection{Comment~2.1 -- Title}

\subsubsection{Reviewer Comment}
\begin{reviewerbox}
\reviewlabel{} I suggest you improve the title of this paper to one that is
much more descriptive and suitable for this research.
\end{reviewerbox}

\subsubsection{Authors' Response}
\begin{responsebox}
\responselabel{}

% TODO: Insert response here.
\end{responsebox}

%% ------------------------------------------------------------
\subsection{Comment~2.2 -- Unassigned affiliation (Line 18)}

\subsubsection{Reviewer Comment}
\begin{reviewerbox}
\reviewlabel{} (Line 18) ``Homi Bhabha National Institute, Training School
Complex, Anushaktinagar, Mumbai 400094, India'' --- this affiliation is not
assigned to an author.
\end{reviewerbox}

\subsubsection{Authors' Response}
\begin{responsebox}
\responselabel{}

% TODO: Insert response here.
\end{responsebox}

%% ------------------------------------------------------------
\subsection{Comment~2.3 -- Distinction between neutral species and molecules (Lines 53--54)}

\subsubsection{Reviewer Comment}
\begin{reviewerbox}
\reviewlabel{} (Lines 53--54) There is no clear difference here between a
neutral species and a molecule. A species can be either neutral or a molecule.
\end{reviewerbox}

\subsubsection{Authors' Response}
\begin{responsebox}
\responselabel{}

% TODO: Insert response here.
\end{responsebox}

%% ------------------------------------------------------------
\subsection{Comment~2.4 -- Additional studies on photolysis and radiolysis of pure oxygen (p.~4, Lines 30--37)}

\subsubsection{Reviewer Comment}
\begin{reviewerbox}
\reviewlabel{} (p.~4, Lines 30--37) There are other studies on the photolysis
and radiolysis of pure oxygen. Could their results be applied to those
presented by Zhen \& Linnartz (2013) or by Jamieson et al.\ (2007)?
\begin{itemize}[nosep]
  \item Zhen \& Linnartz (2013): UV-induced photodesorption and photochemistry
        of O$_2$ ice.
  \item Sivaraman et al.: Temperature-dependent formation of ozone in solid
        oxygen by 5~keV electron irradiation and implications for Solar System
        ices.
\end{itemize}
\end{reviewerbox}

\subsubsection{Authors' Response}
\begin{responsebox}
\responselabel{}

% TODO: Insert response here.
\end{responsebox}

%% ------------------------------------------------------------
\subsection{Comment~2.5 -- Undefined acronym COMs (p.~4, Line 17)}

\subsubsection{Reviewer Comment}
\begin{reviewerbox}
\reviewlabel{} (p.~4, Line 17) The acronym COMs has not been previously
defined.
\end{reviewerbox}

\subsubsection{Authors' Response}
\begin{responsebox}
\responselabel{}

% TODO: Insert response here.
\end{responsebox}

%% ------------------------------------------------------------
\subsection{Comment~2.6 -- Identical reactions in processes P7 and P10 (Tables 1 and 5)}

\subsubsection{Reviewer Comment}
\begin{reviewerbox}
\reviewlabel{} (Tables 1 and 5) Please explain why the chemical reactions in
processes P7 and P10 are the same.
\end{reviewerbox}

\subsubsection{Authors' Response}
\begin{responsebox}
\responselabel{}

% TODO: Insert response here.
\end{responsebox}

%% ------------------------------------------------------------
\subsection{Comment~2.7 -- Entropy of activation sentence (p.~5, Lines 53--56)}

\subsubsection{Reviewer Comment}
\begin{reviewerbox}
\reviewlabel{} (p.~5, Lines 53--56) Please explain the following sentence in
more detail or cite a previous work: ``The underlying causes for such
variation are not, as far as we know, currently understood but could be
related to effects such as the entropy of activation.''
\end{reviewerbox}

\subsubsection{Authors' Response}
\begin{responsebox}
\responselabel{}

% TODO: Insert response here.
\end{responsebox}

%% ------------------------------------------------------------
\subsection{Comment~2.8 -- Discrepancy in O$_3$/O$_2$ percentage (Fig.~1)}

\subsubsection{Reviewer Comment}
\begin{reviewerbox}
\reviewlabel{} (Fig.~1) Gerakines et al.\ (1996) reported a percentage
(O$_3$/O$_2$) of 36\%, but Figure~1 shows a percentage of about 24\%. Please
explain this difference.
\end{reviewerbox}

\subsubsection{Authors' Response}
\begin{responsebox}
\responselabel{}

% TODO: Insert response here.
\end{responsebox}

\bigskip
\hrule
\bigskip

\end{document}